\chapter{Model Theory}\label{model-theory}
\chapterstatus{complete}

\section{Introduction}\label{model-theory:section-introduction}

Chapter~\ref{logic} proved the compactness and L\"owenheim--Skolem theorems for countable languages, as consequences of the completeness theorem, and
without the axiom of choice. With the set theory of Chapter~\ref{sets} (cardinals and choice) we can now say how large the models of a first-order
theory can be. The downward L\"owenheim--Skolem theorem (Theorem~\ref{model-theory:theorem-downward-lowenheim-skolem}) produces small elementary
substructures; ultraproducts and Łoś's theorem (Theorem~\ref{model-theory:theorem-los}) give the compactness theorem for languages of any size
(Theorem~\ref{model-theory:theorem-compactness}); and together they give the upward L\"owenheim--Skolem theorem
(Theorem~\ref{model-theory:theorem-upward-lowenheim-skolem}): a first-order theory with an infinite model has models of every infinite cardinality at
least that of the language. So first-order logic cannot pin down an infinite structure up to isomorphism.

In the other direction, the Łoś--Vaught test (Theorem~\ref{model-theory:theorem-los-vaught}) shows that a theory with a unique model of some
infinite cardinality is complete, and Cantor's back-and-forth argument (Theorem~\ref{model-theory:theorem-cantor}) supplies an example: the dense linear
orders. Back-and-forth arguments are formalised by Ehrenfeucht--Fra\"iss\'e theory (Section~\ref{model-theory:section-ehrenfeucht-fraisse}), which
leads to Lindstr\"om's theorem (Theorem~\ref{model-theory:theorem-lindstrom}): first-order logic is the strongest logic that has both the compactness
and the L\"owenheim--Skolem property. References: \cite{chang-keisler}, \cite{marker-model-theory}, \cite[Chapters XII and XIII]{ebbinghaus-flum-thomas}.

\noindent\textbf{Prerequisites.} Chapter~\ref{logic} (Mathematical Logic), Chapter~\ref{sets} (Set Theory) and Chapter~\ref{number-systems} (Number
Systems).

\begin{notation}\label{model-theory:notation-conventions}
We work in ZFC. Languages, formulas and structures are as in Chapter~\ref{logic}, but a language $\mathcal{L}$ may now have any cardinality; $|\mathcal{L}|$
is the cardinality of its set of constant, function and relation symbols, and $\|\mathcal{L}\| = \max(|\mathcal{L}|, \aleph_0)$. If the free variables of a
formula $\varphi$ are among the distinct variables $x_1, \ldots, x_n$, we write $\varphi(x_1, \ldots, x_n)$, and for elements $a_1, \ldots, a_n$ of a structure
$\mathcal{A}$ we write $\mathcal{A} \models \varphi[a_1, \ldots, a_n]$ if $\mathcal{A} \models \varphi[s]$ for an assignment $s$ with $s(x_i) = a_i$; by the
coincidence lemma (Lemma~\ref{logic:lemma-coincidence}) this does not depend on $s$. Tuples are written $\bar{a} = (a_1, \ldots, a_n)$, and $\mathcal{A}$ also
denotes the universe of $\mathcal{A}$. If $\mathcal{L} \subseteq \mathcal{L}'$ and $\mathcal{A}'$ is an $\mathcal{L}'$-structure, its \emph{reduct}
$\mathcal{A}'|_{\mathcal{L}}$ forgets the symbols not in $\mathcal{L}$, and $\mathcal{A}'$ is an \emph{expansion} of $\mathcal{A}'|_{\mathcal{L}}$; by induction on
formulas, an $\mathcal{L}$-formula has the same truth value in $\mathcal{A}'$ and in $\mathcal{A}'|_{\mathcal{L}}$.
\end{notation}

\begin{lemma}\label{model-theory:lemma-counting-formulas}
A language $\mathcal{L}$ has at most $\|\mathcal{L}\|$ terms and at most $\|\mathcal{L}\|$ formulas.
\end{lemma}

\begin{proof}
Let $\kappa = \|\mathcal{L}\|$. Let $T_0$ be the set of variables and constant symbols, and $T_{k+1}$ the union of $T_k$ with the set of terms
$f(t_1, \ldots, t_n)$ with $t_i \in T_k$. By induction $|T_k| \leq \kappa$: the new terms are indexed by the pairs of a function symbol and a finite tuple
of elements of $T_k$, and there are at most $\kappa \cdot \kappa = \kappa$ of these, by Proposition~\ref{sets:proposition-union-bound} and
Theorem~\ref{sets:theorem-hessenberg}. The union of the $T_k$ contains the variables and constants and is closed under the function symbols, so it contains
every term, and it has cardinality at most $\kappa$ by Proposition~\ref{sets:proposition-union-bound}. The same argument, starting from the atomic formulas
and closing under $\neg$, $\to$ and $\forall x$, bounds the formulas.
\end{proof}

\section{Elementary substructures}\label{model-theory:section-elementary}

\begin{definition}\label{model-theory:definition-elementary}
Let $\mathcal{A}$ and $\mathcal{B}$ be $\mathcal{L}$-structures.
\begin{enumerate}
\item An \emph{embedding} $h \colon \mathcal{A} \to \mathcal{B}$ is an injective map preserving constants, functions and relations as in
Definition~\ref{logic:definition-isomorphism-structures} (relations in both directions), but not necessarily surjective. If $\mathcal{A} \subseteq \mathcal{B}$ as
sets and the inclusion is an embedding, $\mathcal{A}$ is a \emph{substructure} of $\mathcal{B}$.
\item An embedding $h$ is \emph{elementary} if $\mathcal{A} \models \varphi[\bar{a}]$ if and only if $\mathcal{B} \models \varphi[h(\bar{a})]$ for every formula
$\varphi(x_1, \ldots, x_n)$ and every $\bar{a} \in \mathcal{A}^n$. A substructure whose inclusion is elementary is an \emph{elementary substructure}, written
$\mathcal{A} \preceq \mathcal{B}$; then $\mathcal{B}$ is an \emph{elementary extension} of $\mathcal{A}$.
\item $\mathcal{A}$ and $\mathcal{B}$ are \emph{elementarily equivalent}, $\mathcal{A} \equiv \mathcal{B}$, if they satisfy the same sentences, that is,
$\Th(\mathcal{A}) = \Th(\mathcal{B})$.
\end{enumerate}
\end{definition}

A subset of $\mathcal{B}$ is the universe of a substructure if and only if it is nonempty, contains the constants $c^{\mathcal{B}}$ and is closed under the
functions $f^{\mathcal{B}}$; the substructure is then unique. Isomorphisms are elementary (Proposition~\ref{logic:proposition-isomorphism-invariance}), and
$\mathcal{A} \preceq \mathcal{B}$ implies $\mathcal{A} \equiv \mathcal{B}$, by taking $n = 0$. By induction on terms, a substructure computes the values of
terms at elements of $\mathcal{A}$ as $\mathcal{B}$ does, so quantifier-free formulas have the same truth values in $\mathcal{A}$ and $\mathcal{B}$; the
difference is in the quantifiers.

\begin{example}\label{model-theory:example-substructures}
\begin{enumerate}
\item In the language of orders, $(\NN, <)$ is a substructure of $(\ZZ, <)$ but not an elementary one: for $\varphi(x) = \exists y\,(y < x)$,
$(\ZZ, <) \models \varphi[0]$ and $(\NN, <) \not\models \varphi[0]$.
\item $\mathcal{A} \subseteq \mathcal{B}$ and $\mathcal{A} \equiv \mathcal{B}$ do not imply $\mathcal{A} \preceq \mathcal{B}$. The structure $(\NN_{>0}, <)$ is
a substructure of $(\NN, <)$ and is isomorphic to it by $n \mapsto n - 1$, so the two are elementarily equivalent; but $1$ has a predecessor in $\NN$ and none
in $\NN_{>0}$, so the same formula $\varphi$ shows that the inclusion is not elementary.
\item In the language of rings, $\ZZ$ is a substructure of $\QQ$, and not even elementarily equivalent to it: the sentence $\exists y\,(y + y = 1)$ holds in $\QQ$
and not in $\ZZ$.
\end{enumerate}
\end{example}

The following test only involves truth in the larger structure.

\begin{lemma}[Tarski--Vaught test]\label{model-theory:lemma-tarski-vaught}
Let $\mathcal{A}$ be a substructure of $\mathcal{B}$. Then $\mathcal{A} \preceq \mathcal{B}$ if and only if, for every formula $\varphi$, every variable $x$
and every assignment $s$ with values in $\mathcal{A}$ such that $\mathcal{B} \models \exists x\,\varphi[s]$, there is $a \in \mathcal{A}$ with
$\mathcal{B} \models \varphi[s[x \mapsto a]]$.
\end{lemma}

\begin{proof}
If $\mathcal{A} \preceq \mathcal{B}$ and $\mathcal{B} \models \exists x\,\varphi[s]$, then $\mathcal{A} \models \exists x\,\varphi[s]$, so there is $a \in \mathcal{A}$
with $\mathcal{A} \models \varphi[s[x \mapsto a]]$, and then $\mathcal{B} \models \varphi[s[x \mapsto a]]$. Conversely, assume the condition. We show by
induction on $\varphi$ that $\mathcal{A} \models \varphi[s]$ if and only if $\mathcal{B} \models \varphi[s]$ for every assignment $s$ with values in $\mathcal{A}$.
Atomic formulas have the same truth values, as remarked above, and the cases $\neg$ and $\to$ follow from the induction hypothesis. Let
$\varphi = \forall x\,\psi$. By induction, $\mathcal{A} \models \varphi[s]$ if and only if $\mathcal{B} \models \psi[s[x \mapsto a]]$ for all $a \in \mathcal{A}$.
This holds if $\mathcal{B} \models \forall x\,\psi[s]$. If $\mathcal{B} \not\models \forall x\,\psi[s]$, then $\mathcal{B} \models \exists x\,\neg\psi[s]$, and the
condition gives $a \in \mathcal{A}$ with $\mathcal{B} \models \neg\psi[s[x \mapsto a]]$; so both sides fail.
\end{proof}

\begin{example}\label{model-theory:example-rationals-reals}
$(\QQ, <) \preceq (\RR, <)$. We check the Tarski--Vaught test. Let $s$ take values in $\QQ$, and suppose $(\RR, <) \models \varphi[s[x \mapsto b]]$ with
$b \in \RR$. The free variables of $\exists x\,\varphi$ take finitely many values $q_1 < \cdots < q_m$ in $\QQ$. If $b$ is one of them we are done. Otherwise $b$
lies in one of the open intervals into which the $q_j$ cut $\RR$, and this interval contains a rational $r$ (Proposition~\ref{number-systems:proposition-complete-archimedean}).
There is an order-preserving bijection $\sigma \colon \RR \to \RR$ fixing all $q_j$ with $\sigma(b) = r$: on an interval $[q_j, q_{j+1}]$ take the map that is affine on
$[q_j, b]$ and on $[b, q_{j+1}]$, sending $q_j, b, q_{j+1}$ to $q_j, r, q_{j+1}$, and the identity elsewhere; on an unbounded interval use an affine map of the
half line fixing its endpoint, and if $m = 0$ a translation. By Proposition~\ref{logic:proposition-isomorphism-invariance} applied to the automorphism $\sigma$,
$(\RR, <) \models \varphi[s[x \mapsto r]]$, since $\sigma \circ s[x \mapsto b]$ agrees with $s[x \mapsto r]$ on the free variables of $\varphi$. So
$(\QQ, <) \equiv (\RR, <)$, although the two orders are not isomorphic, one being countable and the other not
(Proposition~\ref{number-systems:proposition-reals-cardinality}).
\end{example}

\begin{theorem}[Downward L\"owenheim--Skolem]\label{model-theory:theorem-downward-lowenheim-skolem}
Let $\mathcal{B}$ be an $\mathcal{L}$-structure, $X \subseteq \mathcal{B}$, and $\kappa$ a cardinal with $\|\mathcal{L}\| \leq \kappa$ and
$|X| \leq \kappa \leq |\mathcal{B}|$. Then $\mathcal{B}$ has an elementary substructure $\mathcal{A}$ with $X \subseteq \mathcal{A}$ and $|\mathcal{A}| = \kappa$.
\end{theorem}

\begin{proof}
Since $\kappa \leq |\mathcal{B}|$ there is a subset $Y \subseteq \mathcal{B}$ with $|Y| = \kappa$, and $X_0 = X \cup Y$ has cardinality $\kappa$
(Corollary~\ref{sets:corollary-cardinal-arithmetic}). Consider the triples $(\varphi, x, \bar{b})$ of a formula $\varphi$, a variable $x$ and a tuple $\bar{b}$
of elements of $\mathcal{B}$ whose length is the number of free variables $y_1, \ldots, y_n$ of $\exists x\,\varphi$, listed in the order of their indices, such
that $\mathcal{B} \models \exists x\,\varphi[\bar{b}]$. By the axiom of choice there is a function $w$ choosing for each such triple a \emph{witness}
$w(\varphi, x, \bar{b}) \in \mathcal{B}$ with $\mathcal{B} \models \varphi[s]$ for the assignment sending $x$ to the witness and $y_i$ to $b_i$. Define
$X_{k+1}$ as the union of $X_k$ with the witnesses of all triples whose tuple $\bar{b}$ has entries in $X_k$. There are at most $\kappa$ such triples, by
Lemma~\ref{model-theory:lemma-counting-formulas} and since $X_k$ has at most $\kappa$ finite tuples; so $|X_{k+1}| \leq \kappa$, and
$A = \bigcup_kX_k$ has cardinality $\kappa$.

The set $A$ is closed under the functions of $\mathcal{B}$: for $a_1, \ldots, a_n \in A$, all lying in some $X_k$, the formula $x = f(y_1, \ldots, y_n)$ has the
unique witness $f^{\mathcal{B}}(a_1, \ldots, a_n)$, which lies in $X_{k+1}$. Likewise the witness of $x = c$ is $c^{\mathcal{B}} \in X_1$. So $A$ is the universe
of a substructure $\mathcal{A}$. It satisfies the Tarski--Vaught test: if $s$ takes values in $A$ and $\mathcal{B} \models \exists x\,\varphi[s]$, the finitely
many values of $s$ on the free variables of $\exists x\,\varphi$ lie in some $X_k$, and the corresponding witness lies in $X_{k+1} \subseteq A$. By
Lemma~\ref{model-theory:lemma-tarski-vaught}, $\mathcal{A} \preceq \mathcal{B}$.
\end{proof}

For countable $\mathcal{L}$ and $\kappa = \aleph_0$ this refines Theorem~\ref{logic:theorem-lowenheim-skolem}: every infinite structure has a countable
elementary substructure, which even contains any given countable subset. Some choice is necessary here: the proof chooses witnesses in an arbitrary
structure $\mathcal{B}$, while the proof in Chapter~\ref{logic} built a model from the syntax.

\section{Ultraproducts and compactness}\label{model-theory:section-ultraproducts}

\begin{definition}\label{model-theory:definition-ultrafilter}
A \emph{filter} on a set $I$ is a set $F$ of subsets of $I$ with $I \in F$, $\emptyset \notin F$, and closed under finite intersections and under taking
supersets. It is an \emph{ultrafilter} if moreover, for every $X \subseteq I$, either $X \in F$ or $I \setminus X \in F$. For $i \in I$, the \emph{principal}
ultrafilter at $i$ is $\{X \subseteq I : i \in X\}$. A family of subsets of $I$ has the \emph{finite intersection property} if all its finite intersections are
nonempty.
\end{definition}

\begin{lemma}\label{model-theory:lemma-ultrafilters}
\begin{enumerate}
\item If $U$ is an ultrafilter, then $X \cup Y \in U$ if and only if $X \in U$ or $Y \in U$.
\item Every family of subsets of $I$ with the finite intersection property is contained in an ultrafilter.
\end{enumerate}
\end{lemma}

\begin{proof}
(1) If $X \in U$ then $X \cup Y \in U$. If $X, Y \notin U$, then $I \setminus X$ and $I \setminus Y$ lie in $U$, and so does their intersection, which is disjoint
from $X \cup Y$; so $X \cup Y \notin U$, as $\emptyset \notin U$.

(2) Let $\mathcal{S}$ be the family. The sets containing some finite intersection of members of $\mathcal{S}$ form a filter containing $\mathcal{S}$. Filters
containing $\mathcal{S}$, ordered by inclusion, satisfy the hypothesis of Zorn's lemma (Theorem~\ref{sets:theorem-choice-equivalents}), since the union of a
chain of filters is a filter. Let $U$ be a maximal one and $X \subseteq I$. If $X$ meets every member of $U$, the sets containing some $X \cap Y$ with $Y \in U$
form a filter containing $U$ and $X$, so $X \in U$ by maximality. Otherwise $X \cap Y = \emptyset$ for some $Y \in U$, so $Y \subseteq I \setminus X$ and
$I \setminus X \in U$. So $U$ is an ultrafilter.
\end{proof}

\begin{definition}\label{model-theory:definition-ultraproduct}
Let $(\mathcal{A}_i)_{i \in I}$ be $\mathcal{L}$-structures and $U$ an ultrafilter on $I$. On the product $\prod_i\mathcal{A}_i$ of the universes (nonempty
by the axiom of choice), let $g \sim g'$ if $\{i : g(i) = g'(i)\} \in U$; this is an equivalence relation, transitivity coming from closure under
intersections. The \emph{ultraproduct} $\prod_i\mathcal{A}_i/U$ is the structure with universe the set of classes $[g]$, with
$c = [(c^{\mathcal{A}_i})_i]$, $f([g_1], \ldots, [g_n]) = [(f^{\mathcal{A}_i}(g_1(i), \ldots, g_n(i)))_i]$, and $([g_1], \ldots, [g_n]) \in R$ if and only if
$\{i : (g_1(i), \ldots, g_n(i)) \in R^{\mathcal{A}_i}\} \in U$. These are well defined: if $g_j \sim g'_j$ for all $j$, the set where $g_j(i) = g'_j(i)$ for all
$j$ is in $U$, and on it the two definitions agree. If all $\mathcal{A}_i$ equal $\mathcal{A}$, the ultraproduct is the \emph{ultrapower} $\mathcal{A}^I/U$.
\end{definition}

\begin{theorem}[Łoś]\label{model-theory:theorem-los}
Let $s_i$ be an assignment in $\mathcal{A}_i$ for each $i$, and let $s$ be the assignment in $\prod_i\mathcal{A}_i/U$ with $s(x) = [(s_i(x))_i]$. Then for every
formula $\varphi$,
\[
\textstyle\prod_i\mathcal{A}_i/U \models \varphi[s] \iff \{i \in I : \mathcal{A}_i \models \varphi[s_i]\} \in U.
\]
Every assignment in the ultraproduct is of this form.
\end{theorem}

\begin{proof}
Every assignment is of this form: choose a representative of each value (axiom of choice). Write $\mathcal{P} = \prod_i\mathcal{A}_i/U$ and
$S_\varphi = \{i : \mathcal{A}_i \models \varphi[s_i]\}$. By induction on terms, $t^{\mathcal{P}}[s] = [(t^{\mathcal{A}_i}[s_i])_i]$. For atomic formulas the claim
is then the definition of $\sim$ and of $R^{\mathcal{P}}$. Since $S_{\neg\varphi} = I \setminus S_\varphi$, the case of $\neg$ follows from the ultrafilter property.
Since $S_{\varphi \to \psi} = (I \setminus S_\varphi) \cup S_\psi$, the case of $\to$ follows from Lemma~\ref{model-theory:lemma-ultrafilters}(1).

Let $\varphi = \forall x\,\psi$. For $g \in \prod_i\mathcal{A}_i$, the assignment $s[x \mapsto [g]]$ is obtained from the assignments $s_i[x \mapsto g(i)]$, so by
induction $\mathcal{P} \models \psi[s[x \mapsto [g]]]$ if and only if $T_g = \{i : \mathcal{A}_i \models \psi[s_i[x \mapsto g(i)]]\} \in U$. If $S_\varphi \in U$, then
$T_g \supseteq S_\varphi$ is in $U$ for every $g$, so $\mathcal{P} \models \varphi[s]$. If $S_\varphi \notin U$, choose $g(i)$ with
$\mathcal{A}_i \not\models \psi[s_i[x \mapsto g(i)]]$ for $i \notin S_\varphi$, and arbitrary for $i \in S_\varphi$. Then $T_g \subseteq S_\varphi$, so $T_g \notin U$ and
$\mathcal{P} \not\models \varphi[s]$.
\end{proof}

In particular a structure $\mathcal{A}$ embeds elementarily into each of its ultrapowers, by $a \mapsto [(a)_i]$: for constant families the set in Łoś's
theorem is $I$ or $\emptyset$.

\begin{theorem}[Compactness theorem]\label{model-theory:theorem-compactness}
Let $\mathcal{L}$ be any language and $\Gamma$ a set of $\mathcal{L}$-formulas. If every finite subset of $\Gamma$ is satisfiable, then $\Gamma$ is satisfiable.
\end{theorem}

\begin{proof}
Let $I$ be the set of finite subsets of $\Gamma$, and for $\Delta \in I$ choose a structure $\mathcal{A}_\Delta$ and an assignment $s_\Delta$ satisfying $\Delta$.
The sets $\widehat{\Delta} = \{\Delta' \in I : \Delta \subseteq \Delta'\}$ have the finite intersection property, since
$\widehat{\Delta_1} \cap \cdots \cap \widehat{\Delta_r}$ contains $\Delta_1 \cup \cdots \cup \Delta_r$. Let $U$ be an ultrafilter containing them
(Lemma~\ref{model-theory:lemma-ultrafilters}), and $s$ the assignment in $\prod_\Delta\mathcal{A}_\Delta/U$ obtained from the $s_\Delta$. For $\gamma \in \Gamma$,
the set of $\Delta$ with $\mathcal{A}_\Delta \models \gamma[s_\Delta]$ contains $\widehat{\{\gamma\}} \in U$, so by Theorem~\ref{model-theory:theorem-los} the ultraproduct
satisfies $\gamma$ under $s$.
\end{proof}

For countable languages this is Theorem~\ref{logic:theorem-compactness}, which was deduced from the completeness theorem; the proof above avoids formal
deductions altogether, at the price of the axiom of choice.

\section{Upward L\"owenheim--Skolem and categoricity}\label{model-theory:section-categoricity}

\begin{theorem}[Upward L\"owenheim--Skolem]\label{model-theory:theorem-upward-lowenheim-skolem}
Let $T$ be an $\mathcal{L}$-theory with an infinite model, and $\kappa \geq \|\mathcal{L}\|$ a cardinal. Then $T$ has a model of cardinality exactly $\kappa$.
\end{theorem}

\begin{proof}
Add new constant symbols $c_\alpha$, $\alpha < \kappa$, to get a language $\mathcal{L}'$ with $\|\mathcal{L}'\| = \kappa$ (Corollary~\ref{sets:corollary-cardinal-arithmetic}),
and let $\Gamma = T \cup \{c_\alpha \neq c_\beta : \alpha < \beta < \kappa\}$. A finite subset of $\Gamma$ mentions finitely many new constants; interpreting them
as distinct elements of an infinite model of $T$ (and the other new constants arbitrarily) satisfies it. By Theorem~\ref{model-theory:theorem-compactness},
$\Gamma$ has a model $\mathcal{B}$, and $\alpha \mapsto c_\alpha^{\mathcal{B}}$ is injective, so $|\mathcal{B}| \geq \kappa$. By
Theorem~\ref{model-theory:theorem-downward-lowenheim-skolem} applied to $\mathcal{L}'$, $\mathcal{B}$ has an elementary substructure $\mathcal{A}$ of cardinality
$\kappa$, which satisfies $\Gamma$ like $\mathcal{B}$; its reduct to $\mathcal{L}$ is a model of $T$ of cardinality $\kappa$.
\end{proof}

So a first-order theory with an infinite model never determines it up to isomorphism: it has models of arbitrarily large cardinality. Example~\ref{logic:example-nonstandard-arithmetic}
exhibited a different phenomenon, a non-standard model of the same (countable) cardinality.

\begin{definition}\label{model-theory:definition-categorical}
Let $\kappa$ be an infinite cardinal. A theory $T$ is \emph{$\kappa$-categorical} if it has a model of cardinality $\kappa$ and any two models of cardinality
$\kappa$ are isomorphic. It is \emph{complete} if for every sentence $\sigma$, either $\sigma$ is true in all models of $T$ or $\neg\sigma$ is. For countable
languages this agrees with Definition~\ref{logic:definition-complete-theory}, by the completeness theorem (Corollary~\ref{logic:corollary-completeness}).
\end{definition}

\begin{theorem}[Łoś--Vaught test]\label{model-theory:theorem-los-vaught}
Let $T$ be a satisfiable theory without finite models that is $\kappa$-categorical for some $\kappa \geq \|\mathcal{L}\|$. Then $T$ is complete.
\end{theorem}

\begin{proof}
Suppose $\sigma$ is a sentence such that neither $\sigma$ nor $\neg\sigma$ holds in all models of $T$. Then $T \cup \{\neg\sigma\}$ and $T \cup \{\sigma\}$ have
models, which are infinite, so by Theorem~\ref{model-theory:theorem-upward-lowenheim-skolem} they have models $\mathcal{A}$ and $\mathcal{B}$ of cardinality
$\kappa$. These are models of $T$, hence isomorphic, but $\mathcal{B} \models \sigma$ and $\mathcal{A} \not\models \sigma$, contradicting
Proposition~\ref{logic:proposition-isomorphism-invariance}.
\end{proof}

\begin{definition}\label{model-theory:definition-dlo}
The theory $T_{\mathrm{DLO}}$ of \emph{dense linear orders without endpoints}, in the language of orders, consists of the sentences
$\forall x\,\neg(x < x)$, $\forall x\forall y\forall z\,((x < y \wedge y < z) \to x < z)$, $\forall x\forall y\,(x < y \vee x = y \vee y < x)$,
$\forall x\forall y\,(x < y \to \exists z\,(x < z \wedge z < y))$, $\forall x\exists y\,(y < x)$ and $\forall x\exists y\,(x < y)$.
\end{definition}

Its models are infinite: a model is nonempty and has no largest element. Examples are $(\QQ, <)$, $(\RR, <)$ and every open interval of $\RR$.

\begin{theorem}[Cantor]\label{model-theory:theorem-cantor}
Any two countable dense linear orders without endpoints are isomorphic.
\end{theorem}

\begin{proof}
Let $A$ and $B$ be such orders, enumerated as $a_0, a_1, \ldots$ and $b_0, b_1, \ldots$ (they are countably infinite). \emph{Extension step}: if $p$ is an
order-preserving bijection from a finite subset of $A$ onto a finite subset of $B$ and $a \in A$, there is $b \in B$ such that $p \cup \{(a, b)\}$ is still
order-preserving and injective. Indeed, if $a \in \operatorname{dom}p$ take $b = p(a)$; otherwise the elements of $\operatorname{dom}p$ smaller than $a$ have
images smaller than those of the elements larger than $a$, and we need $b$ strictly between these two finite sets of images, which exists because $B$ is
dense and has no endpoints (and is arbitrary if $p = \emptyset$). The same holds with the roles of $A$ and $B$ exchanged. Starting from $p_0 = \emptyset$, let
$p_{2k+1}$ extend $p_{2k}$ with $a_k$ in its domain, and $p_{2k+2}$ extend $p_{2k+1}$ with $b_k$ in its image (\emph{back and forth}). The union of the
$p_k$ is an order-preserving bijection $A \to B$.
\end{proof}

\begin{corollary}\label{model-theory:corollary-dlo-complete}
$T_{\mathrm{DLO}}$ is $\aleph_0$-categorical and complete. In particular $(\QQ, <) \equiv (\RR, <)$, and a sentence in the language of orders holds in $(\QQ, <)$
if and only if it holds in every dense linear order without endpoints.
\end{corollary}

\begin{proof}
By Theorem~\ref{model-theory:theorem-cantor}, $T_{\mathrm{DLO}}$ is $\aleph_0$-categorical, its language is finite, and it has no finite models; apply
Theorem~\ref{model-theory:theorem-los-vaught} with $\kappa = \aleph_0$.
\end{proof}

\begin{counterexample}\label{model-theory:counterexample-dlo-continuum}
$T_{\mathrm{DLO}}$ is not $2^{\aleph_0}$-categorical. The orders $(\RR, <)$ and $(\RR \setminus \{0\}, <)$ are dense without endpoints and have cardinality
$2^{\aleph_0}$, but they are not isomorphic: in $\RR \setminus \{0\}$ the set of negative numbers is nonempty and bounded above (by $1$) but has no least upper
bound, since every upper bound $u$ is positive and $u/2$ is a smaller one, while in $\RR$ every nonempty subset bounded above has a least upper bound
(Theorem~\ref{number-systems:theorem-reals}); an order isomorphism would preserve this property. By contrast, the theory of algebraically closed fields of a
fixed characteristic is $\kappa$-categorical for every uncountable $\kappa$ (such a field is determined by its transcendence degree, which equals its cardinality),
hence complete by the Łoś--Vaught test; see \cite[Section 2.2]{marker-model-theory} and Chapter~\ref{fields}.
\end{counterexample}

\section{Ehrenfeucht--Fra\"iss\'e theory}\label{model-theory:section-ehrenfeucht-fraisse}

The proof of Cantor's theorem used only one feature of the orders: finite partial isomorphisms can always be extended, forth and back. We now make this
precise and relate it to the quantifier depth of formulas. In this section and the next, languages are \emph{relational}: they have no constant or function
symbols. This is no real restriction, since an $n$-ary function can be replaced by its graph, an $(n+1)$-ary relation, and a constant by a unary relation
holding at one point. In a relational language the terms are the variables.

\begin{definition}\label{model-theory:definition-quantifier-rank}
The \emph{quantifier rank} of a formula is defined by $\mathrm{qr}(\varphi) = 0$ for atomic $\varphi$, $\mathrm{qr}(\neg\varphi) = \mathrm{qr}(\varphi)$,
$\mathrm{qr}(\varphi \to \psi) = \max(\mathrm{qr}(\varphi), \mathrm{qr}(\psi))$ and $\mathrm{qr}(\forall x\,\varphi) = \mathrm{qr}(\varphi) + 1$; so
$\mathrm{qr}(\exists x\,\varphi) = \mathrm{qr}(\varphi) + 1$, and $\wedge$, $\vee$ take the maximum. For structures $\mathcal{A}$, $\mathcal{B}$, tuples
$\bar{a} \in \mathcal{A}^m$, $\bar{b} \in \mathcal{B}^m$ and $n \geq 0$, write $(\mathcal{A}, \bar{a}) \equiv_n (\mathcal{B}, \bar{b})$ if
$\mathcal{A} \models \varphi[\bar{a}] \iff \mathcal{B} \models \varphi[\bar{b}]$ for every formula $\varphi(x_1, \ldots, x_m)$ of quantifier rank at most $n$; for
$m = 0$ write $\mathcal{A} \equiv_n \mathcal{B}$.
\end{definition}

\begin{definition}\label{model-theory:definition-partial-isomorphism}
Let $\mathcal{A}$ and $\mathcal{B}$ be structures for a relational language. A \emph{partial isomorphism} from $\mathcal{A}$ to $\mathcal{B}$ is an injective map
$p$ from a subset of $\mathcal{A}$ to $\mathcal{B}$ with $(a_1, \ldots, a_r) \in R^{\mathcal{A}} \iff (p(a_1), \ldots, p(a_r)) \in R^{\mathcal{B}}$ for every
$r$-ary relation symbol $R$ and all $a_j \in \operatorname{dom}p$. A set $I$ of partial isomorphisms has the \emph{forth property} with respect to a set $J$ if for
every $p \in I$ and $a \in \mathcal{A}$ there is $q \in J$ with $p \subseteq q$ and $a \in \operatorname{dom}q$, and the \emph{back property} if for every
$p \in I$ and $b \in \mathcal{B}$ there is $q \in J$ with $p \subseteq q$ and $b$ in the image of $q$.
\begin{enumerate}
\item An \emph{$n$-back-and-forth system} is a sequence $(I_k)_{0 \leq k \leq n}$ of nonempty sets of partial isomorphisms such that $I_{k+1}$ has the forth and
back properties with respect to $I_k$ for $k < n$. We write $\mathcal{A} \cong_n \mathcal{B}$ if one exists, and
$(\mathcal{A}, \bar{a}) \cong_n (\mathcal{B}, \bar{b})$ if one exists with $\bar{a} \mapsto \bar{b}$ (the map $a_j \mapsto b_j$) in $I_n$.
\item $\mathcal{A}$ and $\mathcal{B}$ are \emph{partially isomorphic}, $\mathcal{A} \cong_p \mathcal{B}$, if there is a nonempty set $I$ of partial isomorphisms
with the forth and back properties with respect to itself.
\end{enumerate}
\end{definition}

For example, the finite order-preserving injections between two dense linear orders without endpoints show that any two such orders are partially
isomorphic: this is the extension step in the proof of Theorem~\ref{model-theory:theorem-cantor}.

\begin{lemma}\label{model-theory:lemma-countable-partial-isomorphism}
Countable partially isomorphic structures are isomorphic.
\end{lemma}

\begin{proof}
Let $I$ be as in the definition, and enumerate $\mathcal{A} = \{a_0, a_1, \ldots\}$ and $\mathcal{B} = \{b_0, b_1, \ldots\}$, with repetitions if the structures
are finite. Choose $p_0 \in I$, then $p_{2k+1} \in I$ extending $p_{2k}$ with $a_k$ in its domain, and $p_{2k+2} \in I$ extending $p_{2k+1}$ with $b_k$ in its
image. The union of the chain $(p_k)$ is a bijection $\mathcal{A} \to \mathcal{B}$ which preserves the relations, since any finitely many elements lie in the
domain of some $p_k$.
\end{proof}

\begin{theorem}[Fra\"iss\'e]\label{model-theory:theorem-fraisse}
Let $(I_k)_{k \leq n}$ be an $n$-back-and-forth system from $\mathcal{A}$ to $\mathcal{B}$, let $k \leq n$, $p \in I_k$, $\varphi$ a formula with
$\mathrm{qr}(\varphi) \leq k$, and $s$ an assignment in $\mathcal{A}$ whose values on the free variables of $\varphi$ lie in $\operatorname{dom}p$. Then
$\mathcal{A} \models \varphi[s]$ if and only if $\mathcal{B} \models \varphi[t]$ for any assignment $t$ with $t(x) = p(s(x))$ for the free variables $x$ of
$\varphi$. Consequently $\mathcal{A} \cong_n \mathcal{B}$ implies $\mathcal{A} \equiv_n \mathcal{B}$, the same holds with tuples, and
$\mathcal{A} \cong_p \mathcal{B}$ implies $\mathcal{A} \equiv \mathcal{B}$.
\end{theorem}

\begin{proof}
Induction on $\varphi$. An atomic formula is $y = z$ or $R(y_1, \ldots, y_r)$ with variables $y_j$, and the claim holds because $p$ is injective and
preserves the relations. The cases $\neg$ and $\to$ follow from the induction hypothesis. Let $\varphi = \forall y\,\psi$, so $k \geq 1$ and
$\mathrm{qr}(\psi) \leq k - 1$. Suppose $\mathcal{A} \models \varphi[s]$, and let $b \in \mathcal{B}$. By the back property there is $q \in I_{k-1}$ with
$p \subseteq q$ and $b = q(a)$ for some $a$. Then $\mathcal{A} \models \psi[s[y \mapsto a]]$, the values of $s[y \mapsto a]$ on the free variables of $\psi$ lie in
$\operatorname{dom}q$, and $q(s[y \mapsto a](x)) = t[y \mapsto b](x)$ for these variables; by induction $\mathcal{B} \models \psi[t[y \mapsto b]]$. As $b$ was
arbitrary, $\mathcal{B} \models \varphi[t]$. The converse uses the forth property in the same way. For sentences the condition on $s$ is empty, which gives
the consequences, taking $I_k = I$ for all $k$ in the last one.
\end{proof}

\begin{example}\label{model-theory:example-empty-language}
Let $\mathcal{L}$ be the empty language, so structures are nonempty sets. If $|\mathcal{A}| = |\mathcal{B}|$, or both have at least $n$ elements, then
$\mathcal{A} \cong_n \mathcal{B}$: let $I_k$ be the set of injections from subsets of $\mathcal{A}$ of size at most $n - k$ into $\mathcal{B}$ (all injections with
finite domain, if $|\mathcal{A}| = |\mathcal{B}|$ is finite); an injection with at most $n - k - 1$ elements in its domain can be extended by any point, since the
other set has more than $n - k - 1$ elements. So a sentence of quantifier rank $n$ cannot distinguish a set with $n$ elements from one with $n + 1$, and no
sentence of the empty language is true exactly in the finite sets of even cardinality.
\end{example}

For a \emph{finite} relational language the converse of Fra\"iss\'e's theorem holds, and each class $\equiv_n$ is defined by one sentence. Fix an enumeration
of all formulas of $\mathcal{L}$ (a countable set) and the variables $x_1, x_2, \ldots$; the conjunction $\bigwedge\Phi$ or disjunction $\bigvee\Phi$ of a finite
nonempty set $\Phi$ of formulas is taken in the order of this enumeration.

\begin{definition}\label{model-theory:definition-hintikka}
Let $\mathcal{L}$ be a finite relational language. For a structure $\mathcal{A}$, a tuple $\bar{a} \in \mathcal{A}^m$ and $n$ with $m + n \geq 1$, the
\emph{Hintikka formula} $\varphi^n_{\mathcal{A}, \bar{a}}(x_1, \ldots, x_m)$ is defined by recursion on $n$. For $m \geq 1$, $\varphi^0_{\mathcal{A}, \bar{a}}$ is
the conjunction of the atomic formulas in the variables $x_1, \ldots, x_m$ that hold at $\bar{a}$ and the negations of those that do not; there are finitely
many, as $\mathcal{L}$ is finite. Then
\[
\varphi^{n+1}_{\mathcal{A}, \bar{a}} = \varphi^0_{\mathcal{A}, \bar{a}} \wedge \bigwedge_{a \in \mathcal{A}}\exists x_{m+1}\,\varphi^n_{\mathcal{A}, \bar{a}a}
\wedge \forall x_{m+1}\bigvee_{a \in \mathcal{A}}\varphi^n_{\mathcal{A}, \bar{a}a},
\]
where the conjunction and the disjunction are over the \emph{set} of formulas $\{\varphi^n_{\mathcal{A}, \bar{a}a} : a \in \mathcal{A}\}$, which is finite by
Lemma~\ref{model-theory:lemma-hintikka}, and the first conjunct is omitted if $m = 0$.
\end{definition}

\begin{lemma}\label{model-theory:lemma-hintikka}
Let $\mathcal{L}$ be a finite relational language and $m + n \geq 1$. The set $H^n_m$ of all formulas $\varphi^n_{\mathcal{A}, \bar{a}}$, for all structures
$\mathcal{A}$ and $\bar{a} \in \mathcal{A}^m$, is finite. Each has quantifier rank $n$, and $\mathcal{A} \models \varphi^n_{\mathcal{A}, \bar{a}}[\bar{a}]$.
\end{lemma}

\begin{proof}
Induction on $n$, for all $m$ at once; this also shows that Definition~\ref{model-theory:definition-hintikka} makes sense. The set $H^0_m$ ($m \geq 1$) is finite,
since $\varphi^0_{\mathcal{A}, \bar{a}}$ is determined by the subset of the finitely many atomic formulas in $x_1, \ldots, x_m$ that hold at $\bar{a}$. The
formula $\varphi^{n+1}_{\mathcal{A}, \bar{a}}$ is determined by $\varphi^0_{\mathcal{A}, \bar{a}}$ and by a subset of the finite set $H^n_{m+1}$, so $H^{n+1}_m$ is
finite. The statements about quantifier rank and truth follow by the same induction.
\end{proof}

\begin{theorem}[Ehrenfeucht--Fra\"iss\'e]\label{model-theory:theorem-ehrenfeucht-fraisse}
Let $\mathcal{L}$ be a finite relational language, $\bar{a} \in \mathcal{A}^m$, $\bar{b} \in \mathcal{B}^m$ and $m + n \geq 1$. The following are equivalent:
\begin{enumerate}
\item[(a)] $\mathcal{B} \models \varphi^n_{\mathcal{A}, \bar{a}}[\bar{b}]$;
\item[(b)] $(\mathcal{A}, \bar{a}) \cong_n (\mathcal{B}, \bar{b})$;
\item[(c)] $(\mathcal{A}, \bar{a}) \equiv_n (\mathcal{B}, \bar{b})$.
\end{enumerate}
In particular, for $n \geq 1$ the class of structures $\mathcal{B}$ with $\mathcal{B} \equiv_n \mathcal{A}$ is the class of models of the single sentence
$\varphi^n_{\mathcal{A}}$, and there are finitely many such classes.
\end{theorem}

\begin{proof}
(c) implies (a) by Lemma~\ref{model-theory:lemma-hintikka}, and (b) implies (c) by Theorem~\ref{model-theory:theorem-fraisse}. Assume (a). For $k \leq n$ let
$I_k$ be the set of maps $\bar{c} \mapsto \bar{d}$, where $\bar{c} \in \mathcal{A}^{m+n-k}$ and $\bar{d} \in \mathcal{B}^{m+n-k}$ extend $\bar{a}$ and $\bar{b}$ and
$\mathcal{B} \models \varphi^k_{\mathcal{A}, \bar{c}}[\bar{d}]$. Such a map is a well defined partial isomorphism: if the tuples are nonempty,
$\varphi^k_{\mathcal{A}, \bar{c}}$ contains the conjunct $\varphi^0_{\mathcal{A}, \bar{c}}$, which records which equalities $x_i = x_j$ and which relations hold at
$\bar{c}$, and these then hold at $\bar{d}$ as well; if they are empty, it is the empty map. By (a), $\bar{a} \mapsto \bar{b}$ lies in $I_n$. Let
$\bar{c} \mapsto \bar{d}$ be in $I_{k+1}$. For $a \in \mathcal{A}$, the conjunct $\exists x\,\varphi^k_{\mathcal{A}, \bar{c}a}$ of
$\varphi^{k+1}_{\mathcal{A}, \bar{c}}$ gives $b$ with $\mathcal{B} \models \varphi^k_{\mathcal{A}, \bar{c}a}[\bar{d}b]$, so $\bar{c}a \mapsto \bar{d}b$ is in $I_k$:
this is the forth property. For $b \in \mathcal{B}$, the conjunct $\forall x\bigvee_a\varphi^k_{\mathcal{A}, \bar{c}a}$ gives $a$ with $\bar{c}a \mapsto \bar{d}b$ in
$I_k$: the back property. The forth property also shows inductively that all $I_k$ are nonempty. So (b) holds. The last statement follows, with
Lemma~\ref{model-theory:lemma-hintikka} for finiteness.
\end{proof}

\section{Lindstr\"om's theorem}\label{model-theory:section-lindstrom}

First-order logic is one of many possible logics: one can allow infinite conjunctions, quantifiers such as ``there are uncountably many'', or quantification
over subsets. Lindstr\"om's theorem says that none of these extensions keeps both the compactness theorem and the L\"owenheim--Skolem theorem.

\begin{definition}\label{model-theory:definition-abstract-logic}
A \emph{logic} $\mathcal{S}$ assigns to each finite relational language $\mathcal{L}$ a set $\mathcal{S}[\mathcal{L}]$, whose elements are called
\emph{$\mathcal{S}$-sentences}, and a relation $\mathcal{A} \models_{\mathcal{S}} \psi$ between $\mathcal{L}$-structures and $\mathcal{S}$-sentences, such that:
\begin{enumerate}
\item (Isomorphism) isomorphic structures satisfy the same $\mathcal{S}$-sentences;
\item (Expansion) if $\mathcal{L} \subseteq \mathcal{L}'$, then $\mathcal{S}[\mathcal{L}] \subseteq \mathcal{S}[\mathcal{L}']$, and an $\mathcal{L}'$-structure
satisfies $\psi \in \mathcal{S}[\mathcal{L}]$ if and only if its reduct to $\mathcal{L}$ does;
\item (Boolean operations) for $\psi, \chi \in \mathcal{S}[\mathcal{L}]$ there are $\neg\psi$ and $\psi \wedge \chi$ in $\mathcal{S}[\mathcal{L}]$ with the usual
meaning;
\item (Relativisation) for $\psi \in \mathcal{S}[\mathcal{L}]$ and a unary relation symbol $P \notin \mathcal{L}$ there is $\psi^P \in \mathcal{S}[\mathcal{L} \cup \{P\}]$
such that an $(\mathcal{L} \cup \{P\})$-structure $\mathcal{A}$ with $P^{\mathcal{A}} \neq \emptyset$ satisfies $\psi^P$ if and only if
$\mathcal{A}^{(P)} \models_{\mathcal{S}} \psi$, where $\mathcal{A}^{(P)}$ is the $\mathcal{L}$-structure with universe $P^{\mathcal{A}}$ and the restrictions of the
relations of $\mathcal{A}$;
\item (First-order logic) for every first-order $\mathcal{L}$-sentence $\sigma$ there is $\sigma' \in \mathcal{S}[\mathcal{L}]$ with the same models.
\end{enumerate}
The logic $\mathcal{S}$ has the \emph{compactness property} if every set of $\mathcal{S}[\mathcal{L}]$-sentences whose finite subsets have models has a model,
and the \emph{L\"owenheim--Skolem property} if every $\mathcal{S}$-sentence with a model has a countable model.
\end{definition}

First-order logic is a logic in this sense, with relativisation defined by replacing each subformula $\forall x\,\varphi$ by $\forall x\,(P(x) \to \varphi)$;
it has both properties, by Theorems~\ref{logic:theorem-compactness} and~\ref{logic:theorem-lowenheim-skolem}. Infinitary logic with countable conjunctions has
the L\"owenheim--Skolem property but not the compactness property, and the logic with the quantifier ``there are uncountably many'' is compact for countable
languages but fails the L\"owenheim--Skolem property \cite[Chapter XIII]{ebbinghaus-flum-thomas}.

\begin{theorem}[Lindstr\"om]\label{model-theory:theorem-lindstrom}
Let $\mathcal{S}$ be a logic with the compactness and the L\"owenheim--Skolem properties. Then for every finite relational language $\mathcal{L}$, every
$\psi \in \mathcal{S}[\mathcal{L}]$ has the same models as some first-order $\mathcal{L}$-sentence.
\end{theorem}

\begin{proof}
\emph{Step 1.} Suppose that $\psi \in \mathcal{S}[\mathcal{L}]$ has no first-order equivalent. We claim that for every $n \geq 1$ there are $\mathcal{L}$-structures
$\mathcal{A}_n \models \psi$ and $\mathcal{B}_n \not\models \psi$ with $\mathcal{A}_n \cong_n \mathcal{B}_n$. Otherwise, for some $n$, every structure $\equiv_n$ to a
model of $\psi$ is a model of $\psi$ (Theorem~\ref{model-theory:theorem-ehrenfeucht-fraisse}). Let $\sigma$ be the disjunction of the finitely many sentences
$\varphi^n_{\mathcal{A}} \in H^n_0$ with $\mathcal{A} \models \psi$ (or $\exists x\,(x \neq x)$ if $\psi$ has no model). A model of $\psi$ satisfies its own Hintikka
sentence, hence $\sigma$; a model of $\sigma$ is $\equiv_n$ to a model of $\psi$, hence a model of $\psi$. So $\psi$ has the same models as $\sigma$, a
contradiction. Fix $n$-back-and-forth systems $(I^n_k)_{k \leq n}$ from $\mathcal{A}_n$ to $\mathcal{B}_n$.

\emph{Step 2: coding.} Let $\mathcal{L}^+ = \mathcal{L} \cup \{P, Q, U, W, C, <, I, G\}$ with $P, Q, U, W, C$ unary, $<$ and $I$ binary and $G$ ternary. Let
$\mathcal{M}_n$ be the $\mathcal{L}^+$-structure whose universe is the disjoint union of $\mathcal{A}_n$, $\mathcal{B}_n$, $\{0, 1, \ldots, n\}$ and
$\bigcup_kI^n_k$, with $P$, $Q$, $U$, $W$ these four parts, $R^{\mathcal{M}_n} = R^{\mathcal{A}_n} \cup R^{\mathcal{B}_n}$ for $R \in \mathcal{L}$, $<$ the order of
$\{0, \ldots, n\}$, $C = \{n\}$, $I = \{(k, w) : w \in I^n_k\}$ and $G = \{(w, a, b) : w(a) = b\}$. Consider the following first-order $\mathcal{L}^+$-sentences:
\begin{enumerate}
\item[(i)] $P$, $Q$, $U$ are nonempty, $W$ is disjoint from $P$, and $C$ holds at exactly one element, which lies in $U$;
\item[(ii)] $<$ is a linear order on $U$ with a least element, and every other element of $U$ has an immediate predecessor in $U$;
\item[(iii)] for every $k$ in $U$ there is $w$ with $I(k, w)$, and $I(k, w)$ implies $U(k)$ and $W(w)$;
\item[(iv)] for every $w$ in $W$, the relation $G(w, x, y)$ is the graph of an injective partial map from $P$ to $Q$ which preserves each relation of
$\mathcal{L}$ in both directions (a finite conjunction, as $\mathcal{L}$ is finite);
\item[(v)] if $k'$ is the immediate predecessor of $k$ in $U$ and $I(k, w)$, then for every $a$ in $P$ there is $w'$ with $I(k', w')$,
$\forall x\forall y\,(G(w, x, y) \to G(w', x, y))$ and $\exists y\,G(w', a, y)$; and symmetrically for every $b$ in $Q$ there is such a $w'$ with
$\exists x\,G(w', x, b)$.
\end{enumerate}
Let $\chi \in \mathcal{S}[\mathcal{L}^+]$ be the conjunction of $\psi^P$, $(\neg\psi)^Q$ and the $\mathcal{S}$-versions of (i)--(v)
(Definition~\ref{model-theory:definition-abstract-logic}(2)--(5)). Then $\mathcal{M}_n \models \chi$: the structures $\mathcal{M}_n^{(P)} = \mathcal{A}_n$ and
$\mathcal{M}_n^{(Q)} = \mathcal{B}_n$ satisfy $\psi$ and $\neg\psi$, and (v) holds because $I^n_k$ has the forth and back properties with respect to $I^n_{k-1}$.

\emph{Step 3: compactness.} For $j \in \NN$ let $\lambda_j$ be the first-order sentence saying that the element of $C$ has at least $j$ predecessors in $U$. A
finite subset of $\{\chi\} \cup \{\lambda_j' : j \in \NN\}$ is satisfied by $\mathcal{M}_n$ for $n$ large, since $n$ has $n$ predecessors. By the compactness
property there is a model $\mathcal{M}$. Call $k \in U^{\mathcal{M}}$ \emph{infinite} if it has infinitely many predecessors; the element of $C^{\mathcal{M}}$ is
infinite. An infinite $k$ is not the least element, so it has an immediate predecessor $k'$ by (ii), and $k'$ is infinite, as every other predecessor of $k$ is a
predecessor of $k'$.

\emph{Step 4: L\"owenheim--Skolem.} Let $J$ be a new unary relation symbol, and expand $\mathcal{M}$ by
$J = \{w : I(k, w) \text{ for some infinite } k\}$. Let $\theta$ be the first-order sentence saying that $J$ is nonempty and contained in $W$, and that for every
$w$ in $J$ and every $a$ in $P$ (respectively $b$ in $Q$) there is $w'$ in $J$ with $G(w, x, y) \to G(w', x, y)$ for all $x, y$ and $a$ in the domain (respectively
$b$ in the image) of $w'$. The expansion satisfies $\chi \wedge \theta'$: $J$ is nonempty by (iii) applied to the element of $C$, and the extensions required
by $\theta$ exist by (v), since the immediate predecessor of an infinite element is infinite. By the L\"owenheim--Skolem property, $\chi \wedge \theta'$ has a
countable model $\mathcal{N}$. Then $\mathcal{N}^{(P)}$ and $\mathcal{N}^{(Q)}$ are countable $\mathcal{L}$-structures with $\mathcal{N}^{(P)} \models \psi$ and
$\mathcal{N}^{(Q)} \not\models \psi$, by relativisation, and by (iv) and $\theta$ the partial maps coded by the elements of $J^{\mathcal{N}}$ form a nonempty set
of partial isomorphisms with the forth and back properties. By Lemma~\ref{model-theory:lemma-countable-partial-isomorphism},
$\mathcal{N}^{(P)} \cong \mathcal{N}^{(Q)}$, which contradicts the isomorphism property of $\mathcal{S}$. So $\psi$ has a first-order equivalent.
\end{proof}

\begin{remark}\label{model-theory:remark-lindstrom}
The theorem is from \cite{lindstrom-1969}; we followed \cite[Chapter XIII]{ebbinghaus-flum-thomas}. It explains why first-order logic is the logic of this
project: its two basic metatheorems, compactness and L\"owenheim--Skolem, characterise it. The price is expressiveness: by
Example~\ref{model-theory:example-empty-language} and Counterexample~\ref{logic:counterexample-finiteness}, finiteness and even cardinality are not first-order
properties, and by Theorem~\ref{model-theory:theorem-upward-lowenheim-skolem} no infinite structure is characterised up to isomorphism.
\end{remark}

\section{Exercises}\label{model-theory:section-exercises}

\begin{exercise}\label{model-theory:exercise-infinitesimals}
Show that on an infinite set $I$ there is an ultrafilter containing all complements of finite sets, and that such an ultrafilter is not principal. For such an
ultrafilter $U$ on $\NN$, show that the ultrapower of the ordered field $(\RR, 0, 1, +, \cdot, <)$ contains a positive element smaller than $1/n$ for every
$n \in \NN_{>0}$, and that it satisfies the same first-order sentences as $\RR$.
\end{exercise}

\begin{solution}
The complements of finite sets have the finite intersection property, since a finite union of finite sets is finite and $I$ is infinite; so they lie in an
ultrafilter $U$ (Lemma~\ref{model-theory:lemma-ultrafilters}). It is not principal, as $I \setminus \{i\} \in U$ for each $i$. The element
$\epsilon = [(1/(k + 1))_k]$ is positive, and for $n \geq 1$ the set $\{k : 1/(k + 1) < 1/n\}$ is cofinite, hence in $U$; by
Theorem~\ref{model-theory:theorem-los}, $0 < \epsilon < 1/n$, where $1/n$ denotes the class of the constant family. The ultrapower is an elementary extension of
$\RR$ by the remark after Theorem~\ref{model-theory:theorem-los}.
\end{solution}

\begin{exercise}\label{model-theory:exercise-dlo-endpoints}
Show that the theory of dense linear orders with a least and a greatest element (and at least two elements) is $\aleph_0$-categorical and complete.
\end{exercise}

\begin{solution}
Such an order is infinite, by density. Given two countable ones, the orders obtained by removing the endpoints are countable dense linear orders without
endpoints, hence isomorphic (Theorem~\ref{model-theory:theorem-cantor}), and the isomorphism extends by sending endpoints to endpoints. The language is finite
and there are no finite models, so Theorem~\ref{model-theory:theorem-los-vaught} applies.
\end{solution}

\begin{exercise}\label{model-theory:exercise-connectedness}
Show that connectedness of graphs is not first-order: there is no set of sentences in the language with one binary relation $E$ whose models are exactly the
connected graphs (symmetric irreflexive $E$ in which any two vertices are joined by a finite path).
\end{exercise}

\begin{solution}
Suppose $T$ were such a set. Add constants $c$, $d$ and, for each $n$, the sentence $\delta_n$ saying that there is no path of length at most $n$ from $c$ to
$d$. Every finite subset of $T \cup \{\delta_n : n\}$ is satisfied by a long path graph with $c$ and $d$ its endpoints, which is connected. By
Theorem~\ref{model-theory:theorem-compactness} there is a model, a connected graph in which $c$ and $d$ are joined by no path, a contradiction.
\end{solution}
